\scanpage{290}
\mspage{- 101 -}\hfill IX
% archive sheet number: 291
% mounting strip and printed label: 419

envisagé dans (2.5.3), et plus précisément, nous verrons que sa connaissance équivaut essentiellement à celle de l'action du groupe de monodromie $I$ sur $U$.

9.2. \uline{Cas} $\handwritten{\ell\ne}p$ . Plaçons nous donc d'abord dans le cas, plus simple, où $\handwritten{\ell\ne}p$, en reprenant alors les notations introduites dans (2.5.3). Notons qu'à torsion près, la relation $V=U^{I}$ signifie aussi que $V^{\handwritten{\perp}}$ est le module engendré par les éléments $gx-x$ ($u\handwritten{\in}U$, $g\handwritten{\in}I$), et qu'on a en tous cas
\[
gx-x\handwritten{\in}V^{\handwritten{\perp}}\quad\text{pour }x\handwritten{\in}U,\ g\handwritten{\in}I\quad .\tag*{(9.2.1)}
\]
On a donc, pour $g\handwritten{\in}I$ :
\[
gx=x+u(g).x\quad ,\tag*{(9.2.2)}
\]
où $g\handwritten{\longmapsto}u(g)$ est un homomorphisme de $U$ dans $V^{\handwritten{\perp}}$ qui s'annule sur $V=U^I$, ou ce qui revient au même
\[
u\handwritten{\in}\operatorname{Hom}(U/V\ ,V^{\handwritten{\perp}})\simeq (U/V)^{\handwritten{\vee}}\otimes V^{\handwritten{\perp}}\quad .\tag*{(9.2.3)}
\]
Dans le cas envisagé de la réduction semi-stable, i.e. quand on a $V^{\handwritten{\perp}}\handwritten{\subset} V$, la formule (9.2.2) exprime que l'opération $g$ dans $U$ est une \uline{transvection} parallèle à $V^{\handwritten{\perp}}=W$, et la composition de ces transvections n'est autre que l'addition des homomorphismes $u$ correspondants (9.2.3). Donc l'opération du groupe d'inertie $I$ sur $U$ s'explicite entièrement alors par un homomorphisme de groupes
\[
I\longrightarrow (U/V)^{\handwritten{\vee}}\otimes V^{\handwritten{\perp}}\quad .\tag*{(9.2.4)}
\]
\margin{\handwritten{$\unsure{\longleftarrow}$}}
